% appendix-assumptions.tex --- Appendix B: Assumption Register
% ASSERT_CONVENTION: natural_units=kB_equals_1, entropy_base=nats

%----------------------------------------------------------------------
\section{Assumption Register}\label{app:assumptions}
%----------------------------------------------------------------------

The main text (Table~\ref{tab:assumptions}) summarizes the seven
framework assumptions entering the master theorem
(Theorem~\ref{thm:master}).  This appendix provides the complete
register, organized into two tiers: \emph{framework assumptions}
(A1--A7), which define the theoretical setting and are common to all
results, and \emph{model-dependent parameters} (A8--A10), which enter
only the order-of-magnitude estimates of
Sec.~\ref{sec:predictions}.  For each assumption we record its
source, severity, current status, what evidence would close or
validate it, and which results depend on it.

%......................................................................
\subsection{Framework assumptions}\label{app:framework-assumptions}
%......................................................................

Table~\ref{tab:assumptions-full} expands on the main-text summary.
Severity levels are: \textsc{standard} (widely accepted, no serious
prospect of failure), \textsc{model choice} (replaceable without
changing qualitative conclusions), \textsc{medium} (plausible but
unproven), and \textsc{high} (strongest assumption in the chain;
failure would weaken key conclusions).

\begin{table*}[t]
  \centering
  \caption{Full framework assumption register (A1--A7).  ``What would
  close it'' describes the minimal result that would either validate
  the assumption or render it unnecessary.  ``Depends'' lists results
  that rely on each assumption.}
  \label{tab:assumptions-full}
  \begin{tabular}{@{}c p{3.5cm} l l l p{4.5cm}@{}}
    \toprule
    ID & Assumption & Source & Severity & Status
       & What would close it \\
    \midrule
    A1 & Finite-dimensional quantum mechanics
       & Paper~5~\cite{Paper5}
       & \textsc{standard}
       & Accepted
       & --- \\[4pt]
    A2 & Thermal contact at temperature~$T$
       & Standard thermodynamics
       & \textsc{standard}
       & Accepted
       & --- \\[4pt]
    A3 & Closed/semi-closed system equilibrates on finite timescales
       & Statistical mechanics
       & \textsc{high}
       & Open
       & Rigorous equilibration bound for self-modeling Hamiltonians
         (e.g., extending~\cite{Short2012} to SWAP lattices) \\[4pt]
    A4 & SWAP lattice dynamics ($H = JF$)
       & Paper~6~\cite{Paper6}
       & \textsc{model choice}
       & Accepted
       & Show qualitative results hold for general nearest-neighbor~$H$ \\[4pt]
    A5 & Experiential density formula:
         $\rhoexp = \MI{B}{M}(1 - \MI{B}{M}/S_B)$
       & Paper~5~\cite{Paper5}
       & \textsc{standard}
       & Accepted
       & --- (definitional within Paper~5 axioms) \\[4pt]
    A6 & Continuum limit of self-modeling lattice yields smooth
         Lorentzian manifold
       & Paper~6~\cite{Paper6}
       & \textsc{medium}
       & Unvalidated
       & Rigorous continuum-limit construction
         (cf.\ causal-set--manifold correspondence) \\[4pt]
    A7 & Emergent spacetime has Lorentzian signature
       & Paper~6 causal structure~\cite{Paper6}
       & \textsc{medium}
       & Unvalidated
       & Derive signature from causal axioms
         (e.g., via Malament's theorem~\cite{Malament1977}) \\
    \bottomrule
  \end{tabular}
\end{table*}

\noindent
Assumption~A3 is the strongest in the chain.  If it fails---e.g., if
recurrence times for self-modeling Hamiltonians are shorter than
expected---then Route~A of the entropy gradient theorem
(Sec.~\ref{sec:egt}) is weakened, though Routes~B and~C survive.
Assumptions~A6 and~A7 are needed only for the SM-observer
complexification requirement [part~(II) of the master theorem]; the
entropy gradient [part~(I)] is independent of both.  A5 (experiential
density) enters via Route~C and part~(II) step~(d).

%......................................................................
\subsection{Model-dependent parameters}\label{app:model-params}
%......................................................................

The order-of-magnitude estimates P7--P9 of Sec.~\ref{sec:predictions}
require three cosmological or biological inputs external to the
framework.  Table~\ref{tab:params} records each parameter, its
fiducial value, plausible range, and sensitivity.

\begin{table*}[t]
  \centering
  \caption{Model-dependent parameters (A8--A10).  These enter only the
  order-of-magnitude predictions P7--P9; the framework predictions
  P1--P6 are independent of all three.  ``Sensitivity'' indicates how
  the Landauer deficit $\Delta S_{\mathrm{Landauer}}$ or the
  exhaustion timescale $\tauex$ scales with each parameter.}
  \label{tab:params}
  \begin{tabular}{@{}c l l l l@{}}
    \toprule
    ID & Parameter & Fiducial value & Plausible range & Sensitivity \\
    \midrule
    A8 & Cosmological entropy budget $\Smax$
       & ${\sim}\,10^{122}$ nats (Bekenstein)
       & $10^{120}$--$10^{124}$
       & $\tauex \propto \Smax$;
         linear in $\Delta S$ \\[4pt]
    A9 & Neural self-modeling dimension $d_M$
       & ${\sim}\,10^{11}$
       & $10^{9}$--$10^{13}$
       & Logarithmic in $\MI{B}{M}$;
         linear in $\Delta S_{\mathrm{Landauer}}$ \\[4pt]
    A10 & Self-modeling update rate $\nu$
        & ${\sim}\,10$ Hz
        & 1--100 Hz
        & Linear in $N$;
          linear in $\tauex$ \\
    \bottomrule
  \end{tabular}
\end{table*}

\noindent
Varying all three parameters simultaneously across their full
plausible ranges shifts $\Delta S_{\mathrm{Landauer}}$ by at most
${\sim}\,6$ orders of magnitude (from ${\sim}\,10^{25}$ to
${\sim}\,10^{31}$~nats).  The 94-order gap with the Penrose deficit
(Eq.~\eqref{eq:gap-94}) narrows to at worst ${\sim}\,88$ orders of
magnitude: the qualitative conclusion is robust.

%......................................................................
\subsection{Dependency map}\label{app:dep-map}
%......................................................................

Table~\ref{tab:dep-map} records which assumptions are required by
each theorem, proposition, and named result in the paper.  An entry
``$+$A3'' indicates that the assumption strengthens the result
(e.g., extends it from a single-cycle bound to a cosmological
statement) but is not required for the core claim.

\begin{table}[t]
  \centering
  \caption{Dependency map: assumptions required by each result.
  Parenthesized entries are needed only for specific sub-results
  (e.g., a particular route or corollary).}
  \label{tab:dep-map}
  \begin{tabular}{@{}l l@{}}
    \toprule
    Result & Assumptions \\
    \midrule
    Landauer bound (Thm.~\ref{thm:landauer})
      & A1, A2 \\
    Coherence loophole closure
      & A1, A2 (+Markovian bath for R2) \\
    Chain theorem (Prop.~\ref{prop:chain})
      & A1, A2, A3 \\
    Entropy gradient, Route~A
      & A1, A2, A3 \\
    Entropy gradient, Route~B
      & A1, A3 \\
    Entropy gradient, Route~C
      & A1, A5 \\
    Three-consequence theorem (Thm.~\ref{thm:three-consequence})
      & A6, A7 \\
    Gap~C resolution (Thm.~\ref{thm:complexification})
      & A1--A7 \\
    Master theorem (Thm.~\ref{thm:master})
      & A1--A7 \\
    Quantitative estimates P7--P9
      & A1--A7 $+$ A8--A10 \\
    \bottomrule
  \end{tabular}
\end{table}

\noindent
The dependency structure has two notable features.  First, the entropy
gradient theorem has three independent routes with different assumption
subsets (Table~\ref{tab:routes}), so no single assumption failure
destroys the conclusion.  Second, the model-dependent parameters
A8--A10 enter only the quantitative estimates; all structural results
(Landauer bound, chain theorem, Gap~C resolution) are independent of
cosmological inputs.
